\documentclass[11pt]{article}
\usepackage[T1]{fontenc}
\usepackage[margin=1in]{geometry}
\usepackage{amsmath,amssymb,amsthm,booktabs,array}
\usepackage{hyperref}
\hypersetup{colorlinks=true,urlcolor=blue,citecolor=blue,linkcolor=blue,
 pdftitle={Extending Orrick's rectangle count to Barker's king-strip conjectures},
 pdfauthor={Haobo Ma; Wenlin Zhang}}
\newtheorem{theorem}{Theorem}
\newtheorem{corollary}[theorem]{Corollary}
\newcommand{\N}{\mathbb{N}}
\title{Extending Orrick's rectangle count\\to Barker's king-strip conjectures}
\author{Haobo Ma\\\small ChronoAI Pte Ltd, Singapore
\and Wenlin Zhang\\\small National University of Singapore\\\small ChronoAI Pte Ltd, Singapore}
\date{November 16, 2026}
\begin{document}
\maketitle
\begin{abstract}
Extending Orrick's rectangle count (2017, answer 2088242, revision 5)
\cite{orrick}, we prove Colin Barker's conjectures in OEIS A233163,
revision 8, FORMULA lines 8--10 of October 9, 2018, by counting
complementary colour pairs and their normalized eight-colour realizations.
More generally, for every $n,k\geq2$ the corresponding rectangle count is
$2^{nk-4}(2^{n-2}+2^{k-2}-1)$.
Orrick's four-colour classification and Britnell--Wildon's path enumeration
(2015, Proposition 8.1) \cite{britnell-wildon} give all seven empirical
column recurrences in the linked OEIS A233168.
They also give Barker's closed-form and generating-function conjectures
in the path entry A233162 and the width-four through width-seven entries
A233164--A233167.
\end{abstract}

\section{The counting problem}

Hardin's entry A233163 defines arrays of height $n$ and width three;
Barker conjectures their full generating function and a closed formula
\cite{hardin,barker-source}. We derive both formulas from known
four-colour counts by an exact normalization argument.
The linked rectangle table A233168 records
empirical recurrences at widths one through seven \cite{rectangles}.
Chang and Shrock's chromatic polynomials give the labelled four-colour
counts of free strips of widths two, three and four \cite{chang-shrock}.
Britnell and Wildon count paired-colour paths \cite[Proposition 8.1]{britnell-wildon}.
Orrick gives the labelled four-colour count for every rectangular board
\cite{orrick}; Eppstein later gives a zone construction \cite{eppstein}.
The additional ingredient here is an exact first-occurrence normalization
with two choices at every later occurrence of a complementary colour pair.
Together these give the literal eight-colour count at every width, including
the missing-colour path cases, and settle all six Barker entries A233162--A233167.

For positive integers $n,k$, let $T(n,k)$ count arrays
$c=(c_{i,j})\in\{0,1,\ldots,7\}^{n\times k}$ subject to these conditions:
\begin{itemize}
\item Two distinct cells with row and column distances at most one must
have unequal values whose sum is not seven. Thus horizontal, vertical,
and both diagonal neighbours are constrained, with no wraparound.
\item The top-left value is zero.
\item In row-major order, every occurrence of $2,3,4,5,6$ has an earlier $1$;
every occurrence of $3,4,5$ has an earlier $2$; and every $4$ has an earlier $3$.
\end{itemize}
Here ``earlier'' means a smaller value of $k(i-1)+(j-1)$, with
$1\leq i\leq n$ and $1\leq j\leq k$. Colours may be absent.
There is no precedence condition on colour $7$.
These are literal arrays, not counts obtained by dividing by a group order.
Write $a(n)=T(n,3)$ for $n\geq1$, and set $a(0)=0$ for the generating series.

\begin{theorem}\label{thm:rectangle}
For all integers $n,k\geq2$,
\begin{equation}\label{eq:rectangle}
 T(n,k)=2^{nk-4}\bigl(2^{n-2}+2^{k-2}-1\bigr).
\end{equation}
\end{theorem}

\section{Complementary pairs and rectangle counts}

Assign colour $u$ the pair label $\beta(u)=\min(u,7-u)$.
The four pairs are $\{0,7\},\{1,6\},\{2,5\},\{3,4\}$.
For $u,v\in\{0,\ldots,7\}$,
\[
 u\ne v\ \hbox{and}\ u+v\ne7
 \quad\Longleftrightarrow\quad \beta(u)\ne\beta(v).
\]
Every $2\times2$ square therefore contains all four pair labels.

\paragraph{Normalization and the colour choices.}
Partition the cells according to their pair labels, temporarily forgetting
the names of the four classes. Order those classes by first occurrence
in row-major order and name them $0,1,2,3$.
The precedence conditions force the first occurrence in these classes
to have colour $0,1,2,3$, respectively. Indeed, $6$ requires an earlier $1$,
$5$ requires an earlier $2$, and $4$ requires an earlier $3$; the first
appearances of the latter three pairs must occur in that order.
The first occurrence of pair $0$ is the top-left cell and is fixed at $0$.
Every later occurrence of a pair has either of its two colours available:
all required earlier colours have already appeared. Conversely, these
choices satisfy every precedence clause.

Since $n,k\geq2$, all four pairs occur. Each pair partition thus has
exactly $2^{nk-4}$ literal colour realizations. It remains to count the
pair partitions. This is a bijection with independent binary choices,
not an orbit-divisor assumption.

\paragraph{Counting the pair partitions (Orrick's classification).}
We include a proof of Orrick's rectangle count \cite{orrick}, expressed
in terms of horizontal edge sets.
For this count, label the upper-left square as
$\left(\begin{smallmatrix}0&1\\2&3\end{smallmatrix}\right)$.
Every partition has a unique such auxiliary labeling; relabeling by first
appearances recovers the normalization just described.
In the first two rows, each new column must contain the two labels
complementary to the preceding column's labels, in either order.
There are therefore $2^{k-2}$ possible two-row patterns.

For the top row $p_{1,1},\ldots,p_{1,k}$ of such a pattern, put
$H_j=\{p_{1,j},p_{1,j+1}\}$ for $1\leq j<k$.
The four-label condition on each square implies that the horizontal
edge sets of successive rows alternate between $H_j$ and its complement
in $\{0,1,2,3\}$.

If all $H_j$ are equal, the top row alternates between $0$ and $1$,
and the second row alternates between $2$ and $3$.
There is exactly one such initial pattern. Each additional row alternates
between the complementary two labels, with two possible starting labels.
This case contributes $2^{n-2}$ partitions.

Otherwise, some adjacent $H_j,H_{j+1}$ differ. They intersect in
exactly one label, as do their complements. In every row, the shared cell
at column $j+1$ is consequently forced. Its value and the prescribed
two-element edge sets determine the entire row, proceeding in both
directions. A continuation exists: repeat the initial two rows
alternately. Thus each of the other $2^{k-2}-1$ initial patterns has
exactly one continuation to height $n$.
The number of partitions is
$2^{n-2}+2^{k-2}-1$, proving Theorem~\ref{thm:rectangle}.

\begin{corollary}\label{cor:labelled}
For $n,k\geq2$, the number of proper colourings of the king rectangle
with four named colours is $24(2^{n-2}+2^{k-2}-1)$.
Moreover $T(n,k)=T(k,n)$, and
$T(m,m)=2^{m^2-4}(2^{m-1}-1)$ for $m\geq2$.
\end{corollary}
\begin{proof}
Every pair partition uses all four labels, so it has exactly $4!$
labelings. The labelled formula is Orrick's \cite{orrick};
the remaining identities follow from Theorem~\ref{thm:rectangle}.
For square boards the labelled count is $12\cdot2^m-24$,
also given by Eppstein \cite{eppstein}.
\end{proof}

\section{Barker's conjectures and the linked recurrences}

\begin{corollary}\label{cor:barker}
The three-column counts satisfy
\begin{align}
 a(1)&=3, & a(n)&=8^{n-2}(2^n+4)\quad(n\geq2),\label{eq:closed}\\
 a(n)&=24a(n-1)-128a(n-2)\quad(n\geq4),\label{eq:recurrence}\\
 \sum_{n\geq0}a(n)x^n
 &=\frac{x(3-64x+288x^2)}{(1-8x)(1-16x)}
 \quad\hbox{in }\mathbb{Q}[[x]].\label{eq:gf}
\end{align}
\end{corollary}
\begin{proof}
The one-row arrays are $010,012,017$, so $a(1)=3$.
Putting $k=3$ in \eqref{eq:rectangle} gives
$a(n)=4\cdot16^{n-2}+4\cdot8^{n-2}$ for $n\geq2$,
which equals \eqref{eq:closed}. Both geometric terms satisfy
\eqref{eq:recurrence}. Summing them from $n=2$ yields
\[
 3x+\frac{4x^2}{1-16x}+\frac{4x^2}{1-8x}
 =\frac{3x-64x^2+288x^3}{(1-8x)(1-16x)}.
\]
The denominators have constant coefficient one, so these are valid formal
power-series identities. In particular $a(0)=0$, $a(2)=8$, and $a(3)=96$.
\end{proof}

\begin{corollary}\label{cor:recurrences}
For $k\geq3$ and $n\geq4$,
\begin{equation}\label{eq:widthrec}
 T(n,k)=3\cdot2^kT(n-1,k)-2^{2k+1}T(n-2,k).
\end{equation}
For $n\geq2$, $T(n,2)=8^{n-2}$.
\end{corollary}
\begin{proof}
Formula \eqref{eq:rectangle}, viewed as a function of $n$,
is a linear combination of $(2^k)^n$ and $(2^{k+1})^n$.
Their sum and product give the coefficients in \eqref{eq:widthrec}.
At width two the first term disappears from that linear combination,
giving $T(n,2)=8^{n-2}$ for $n\geq2$, and hence
$T(n,2)=8T(n-1,2)$ for $n\geq3$.
This also proves the shift to the powers of eight, A001018, stated in
A233168's cross-references.
\end{proof}

The width-one case is a path rather than a rectangle containing a square.
Let $L(m)=T(m,1)=T(1,m)$.
Let $S(r,s)$ be the number of partitions of an $r$-element set into
$s$ nonempty blocks, with $S(0,0)=1$ and $S(r,s)=0$ for impossible indices.
A path of $m$ cells has $S(m-1,j-1)$ proper pair partitions using $j$
classes. To see this, appending a cell to a partition with $j$ classes
allows any of $j-1$ existing classes other than the preceding one,
or creates a new class from a partition with $j-1$ classes.
This is the recurrence for $S(m-1,j-1)$, with the same initial conditions.
\begin{corollary}\label{cor:path}
For every $m\geq1$ (Britnell--Wildon's path count),
\begin{equation}\label{eq:path}
 L(m)=\sum_{j=1}^{\min(4,m)}2^{m-j}S(m-1,j-1).
\end{equation}
For $m\geq5$, $L(m)=12L(m-1)-44L(m-2)+48L(m-3)$.
For $m\geq2$ and as a formal power series, respectively,
\begin{align}
L(m)&=\frac{2^{m-6}}{9}(90+9\cdot2^m+2\cdot3^m),\label{eq:pathclosed}\\
\sum_{m\geq1}L(m)x^m
 &=\frac{x(1-11x+35x^2-29x^3)}{(1-2x)(1-4x)(1-6x)}.\label{eq:pathgf}
\end{align}
Thus Barker's two conjectures in A233162 are consequences of the path
enumeration \cite{path-source,britnell-wildon}.
\end{corollary}
\begin{proof}
The colour-choice argument above, now with $j$ first occurrences fixed,
gives \eqref{eq:path}.
For $m\geq2$, expand the Stirling numbers using
$S(r,1)=1$, $S(r,2)=(2^r-2)/2$, and
$S(r,3)=(3^r-3\cdot2^r+3)/6$ for $r\geq1$.
Equation~\eqref{eq:path} becomes
\[
 L(m)=\frac{5}{32}2^m+\frac{1}{64}4^m+\frac{1}{288}6^m.
\]
The roots $2,4,6$ give, for $m\geq5$,
\[
 L(m)=12L(m-1)-44L(m-2)+48L(m-3).
\]
The initial values are $L(1),L(2),L(3),L(4)=1,1,3,11$.
The same ordered path represents a horizontal or vertical array.
The exponential expression is \eqref{eq:pathclosed}. Summing from $m=2$
and inserting $L(1)=1$ gives
\[
x+x^2\left(\frac{5/8}{1-2x}+\frac{1/4}{1-4x}
+\frac{1/8}{1-6x}\right),
\]
which has the numerator in \eqref{eq:pathgf}.
Britnell--Wildon's Proposition 8.1 counts these same canonical
paired-colour paths by marked partitions; the proof here uses the
equivalent weighted ordinary-partition formula.
\end{proof}

\begin{corollary}\label{cor:allseries}
For $k\geq2$, put $C_k=2^{k-2}-1$. In $\mathbb Q[[x]]$,
\begin{equation}\label{eq:allseries}
 \sum_{n\geq1}T(n,k)x^n
 =L(k)x+2^{2k-4}x^2\left(
 \frac{1}{1-2^{k+1}x}+\frac{C_k}{1-2^kx}\right).
\end{equation}
This proves all the generating-function and closed-form conjectures
in A233163--A233167 \cite{barker-source,siblings}.
\end{corollary}
\begin{proof}
For $n\geq2$, Theorem~\ref{thm:rectangle} is
$T(n,k)=(2^{k+1})^n/64+C_k(2^k)^n/16$.
Summing these geometric series from $n=2$, and supplying the one-row
coefficient from Corollary~\ref{cor:path}, gives \eqref{eq:allseries}.
The table below expands its numerator over
$(1-2^kx)(1-2^{k+1}x)$. The initial values
$L(3),\ldots,L(7)=3,11,48,236,1248$ follow from \eqref{eq:path}.
The closed forms in the middle column are algebraic rearrangements of
\eqref{eq:rectangle}, valid for $n\geq2$.
\end{proof}
\begin{center}
\small
\begin{tabular}{@{}cll@{}}
\toprule
$k$ & $T(n,k)$, $n\geq2$ & Series numerator\\
\midrule
3 & $8^{n-2}(2^n+4)$ & $x(3-64x+288x^2)$\\
4 & $4^{2n-3}(2^n+12)$ & $x(11-464x+3840x^2)$\\
5 & $2^{5n-6}(2^n+28)$ & $16x(3-256x+4224x^2)$\\
6 & $64^{n-1}(2^n+60)$ & $4x(59-10304x+356352x^2)$\\
7 & $2^{7n-6}(2^n+124)$ & $32x(39-13952x+1019904x^2)$\\
\bottomrule
\end{tabular}
\end{center}

\paragraph{Coverage of the source statements.}
The following maps the FORMULA assertions in all six counting entries
to their proofs. Locations refer to the pinned OEIS exports in the
references. The A233168 table values are instances of the rectangle
and path formulas, rather than additional conjectures.
\begin{center}
\small
\begin{tabular}{@{}>{\raggedright\arraybackslash}p{0.47\linewidth}>{\raggedright\arraybackslash}p{0.46\linewidth}@{}}
\toprule
Source item & Coverage\\
\midrule
A233163 rev.\ 8, lines 7, 9--10 & Proved: Corollary~\ref{cor:barker}\\
A233162 rev.\ 21, lines 9, 11--12 & Proved from prior path count: Corollary~\ref{cor:path}\\
A233164 rev.\ 8, lines 7, 9--10 & Proved: Corollaries~\ref{cor:recurrences}, \ref{cor:allseries}, $k=4$\\
A233165 rev.\ 8, lines 7, 9--10 & Proved: same, $k=5$\\
A233166 rev.\ 8, lines 7, 9--10 & Proved: same, $k=6$\\
A233167 rev.\ 14, lines 7, 9--10 & Proved: same, $k=7$\\
A233168 rev.\ 8, lines 19--25 & Proved: all seven recurrences, Corollaries~\ref{cor:recurrences}, \ref{cor:path}\\
A233168 rev.\ 8, line 31; A001018 & Proved: $T(n,2)=8^{n-2}$, $n\geq2$\\
Chang--Shrock, labelled widths 2--4 & Reproved; prior proof Chang--Shrock, 2000: Corollary~\ref{cor:labelled}\\
Britnell--Wildon, paired-colour paths & Reproved; prior proof Britnell--Wildon, 2015: Corollary~\ref{cor:path}\\
Orrick, labelled rectangle count & Reproved; prior proof Orrick, 2017: Corollary~\ref{cor:labelled}\\
Literal all-width formula and symmetry & Not posed in the OEIS conjectures; proved from prior counts and normalization\\
\bottomrule
\end{tabular}
\end{center}

\section{Related work and verification}

Chang and Shrock compute chromatic polynomials for strips with both
diagonals and free boundaries \cite[Section II]{chang-shrock}.
With four named colours, their length convention is $n=m+2$ vertices.
If $P_k(n)$ denotes this labelled count, their equations (2.1), (2.10)
and (2.14)--(2.22) specialize to
\[
 \sum_{m\geq0}P_k(m+2)z^m=
 \begin{cases}
 24/(1-2z),&k=2,\\
 24(2-3z)/((1-z)(1-2z)),&k=3,\\
 24(4-7z)/((1-z)(1-2z)),&k=4.
 \end{cases}
\]
These are exactly $P_k(n)=24(2^{n-2}+2^{k-2}-1)$ for those widths.
Hardin's A068293, recorded in 2002, also gives one quarter of the square
count \cite{square-data}; that numerical formula predates Barker's
normalized eight-colour conjectures.

Britnell and Wildon prove the paired-colour path enumeration in
Proposition 8.1, with the relevant marked-partition generating functions
in Section 6.5 \cite{britnell-wildon}.
Orrick's answer, first posted January 8, 2017 and expanded in revision 5
on January 9, proves the rectangle count by the two-colour-row or
two-colour-column classification \cite{orrick}.
Later proofs include Blatter's Klein-four-group count of a $4\times5$
board (2017), Eppstein's zone count (2019), paw88789's $4\times4$ count
(2023), anankElpis's square-board proof (2024), and Efremov's published
square-board solution (2025) \cite{later-counts,eppstein,efremov}.

Eppstein counts proper four-colourings of an $m\times m$ king board as
$12\cdot2^m-24$ for $m\geq2$, using two families of crossing zones
\cite[final two paragraphs]{eppstein}. Our partition classification
expresses the same mechanism in terms of adjacent horizontal edge sets.
The normalization fibres and one-row calculation supply the additional
information needed for Barker's eight-colour statements.

At the pinned
\href{https://github.com/the-omega-institute/trureturing/blob/0a197c5f83db6039a3c1177b20d31df4bfe95cd9/D5/S3/Combinatorics/TraceAvoidingArrays/HardinA233163KingStrip.lean}{\texttt{HardinA233163KingStrip.lean}},
the relevant namespace is
\begin{center}\small\ttfamily
D5.S3.Combinatorics.TraceAvoidingArrays.\\
HardinA233163KingStrip
\end{center}
The declaration \texttt{a} defines the literal three-column count and
\texttt{result} proves the complete rational series \eqref{eq:gf}.
Its proof combines the normalized two-row counts, the $16$ and $8$
continuation rates, and an induction over every height.
The all-width formula \eqref{eq:rectangle} and the path argument
\eqref{eq:path} have written proofs here; the linked declaration's
scope is width three.
The imported
\href{https://github.com/the-omega-institute/trureturing/blob/0a197c5f83db6039a3c1177b20d31df4bfe95cd9/D5/S3/Combinatorics/TraceAvoidingArrays/HardinA233196Strip.lean}{\texttt{HardinA233196Strip.lean}}
supplies \texttt{width\_counting\_bridge}, equating the normalized literal
cardinality with a row-transition sum for a fixed positive width.
It is a counting bridge, rather than the rectangle formula.

The accompanying \texttt{verify.py} uses exact integer and rational
arithmetic. It independently counts literal eight-colour arrays and
canonical four-pair arrays over the finite ranges stated in
\texttt{verification.json}. It also checks the two-row seeds, boundary
continuations, recurrence coefficients, and generating-series numerator.
The finite calculations support, rather than replace, the all-size proofs.
The source package contains this TeX file, the checker, its output,
scientific metadata, and rebuild instructions.

\begin{thebibliography}{99}\small\raggedright
\setlength{\itemsep}{0pt}
\bibitem{hardin}
R. H. Hardin, \emph{A233163}, OEIS, revision 8 (October 9, 2018),
NAME and FORMULA.
\url{https://oeis.org/A233163}.
\bibitem{barker-source}
C. Barker, conjectures of October 9, 2018 in Hardin's \emph{A233163},
FORMULA block, revision 8, lines 8--10 of the
\href{https://raw.githubusercontent.com/oeis/oeisdata/571197984aeda7b8b680fe3a1f6a43f9898cd441/seq/A233/A233163.seq}{pinned OEIS export}.
\bibitem{rectangles}
R. H. Hardin, \emph{A233168}, OEIS, revision 8 (October 4, 2024),
FORMULA block, widths one through seven.
\url{https://oeis.org/A233168}.
\bibitem{siblings}
R. H. Hardin, with conjectures by C. Barker (October 9, 2018),
\emph{A233164, A233165, A233166, A233167},
The On-Line Encyclopedia of Integer Sequences.
FORMULA lines 7--10 in revisions 8, 8, 8, 14, respectively;
the last revision is dated September 15, 2024.
\href{https://raw.githubusercontent.com/oeis/oeisdata/571197984aeda7b8b680fe3a1f6a43f9898cd441/seq/A233/A233164.seq}{A233164},
\href{https://raw.githubusercontent.com/oeis/oeisdata/571197984aeda7b8b680fe3a1f6a43f9898cd441/seq/A233/A233165.seq}{A233165},
\href{https://raw.githubusercontent.com/oeis/oeisdata/571197984aeda7b8b680fe3a1f6a43f9898cd441/seq/A233/A233166.seq}{A233166},
\href{https://raw.githubusercontent.com/oeis/oeisdata/571197984aeda7b8b680fe3a1f6a43f9898cd441/seq/A233/A233167.seq}{A233167}.
\bibitem{eppstein}
D. Eppstein, \emph{Coloring kinggraphs}, April 11, 2019.
\url{https://11011110.github.io/blog/2019/04/11/coloring-kinggraphs.html}.
\bibitem{orrick}
W. Orrick, answer 2088242 to \emph{In how many ways can you colorize the
vertices of a grid with 4 colors so that every unit-square vertices are all
of a different color?}, Mathematics Stack Exchange, January 8, 2017;
revision 5, January 9, 2017.
\url{https://math.stackexchange.com/revisions/2088242/5}.
\bibitem{britnell-wildon}
J. R. Britnell and M. Wildon, \emph{Bell numbers, partition moves and the
eigenvalues of the random-to-top shuffle in Dynkin Types A, B and D},
Journal of Combinatorial Theory, Series A \textbf{148} (2017), 116--144.
\href{https://doi.org/10.1016/j.jcta.2016.12.003}{doi:10.1016/j.jcta.2016.12.003}.
\href{https://arxiv.org/abs/1507.04803v1}{arXiv:1507.04803v1},
July 16, 2015, Section 6.5 and Proposition 8.1 (pp.\ 23--24).
\bibitem{path-source}
R. H. Hardin, with conjectures by C. Barker (February 18, 2018),
\emph{A233162}, The On-Line Encyclopedia of Integer Sequences,
revision 21, May 31, 2024, FORMULA lines 9--12 in the
\href{https://raw.githubusercontent.com/oeis/oeisdata/571197984aeda7b8b680fe3a1f6a43f9898cd441/seq/A233/A233162.seq}{pinned export}.
\bibitem{later-counts}
C. Blatter, Mathematics Stack Exchange answer 2398953, August 19, 2017,
\url{https://math.stackexchange.com/a/2398953};
paw88789, answer 4736308, July 14, 2023,
\url{https://math.stackexchange.com/a/4736308};
anankElpis, answer 4980383, October 5, 2024,
\url{https://math.stackexchange.com/a/4980383}.
\bibitem{efremov}
I. Efremov, solution to problem M2810, \emph{Kvant}, 2025, no.\ 1,
pp.\ 23--24. Problem: 2024, no.\ 9, p.\ 20.
\url{https://www.kvant.digital/problems/m2810/}.
\bibitem{chang-shrock}
S.-C. Chang and R. Shrock,
\emph{Ground-state entropy of the Potts antiferromagnet with
next-nearest-neighbor spin-spin couplings on strips of the square lattice},
Phys. Rev. E \textbf{62} (2000), 4650--4664.
\href{https://doi.org/10.1103/PhysRevE.62.4650}{doi:10.1103/PhysRevE.62.4650}.
The formulas cited above are in
\href{https://arxiv.org/abs/cond-mat/0005236v1}{arXiv:cond-mat/0005236v1},
May 15, 2000.
\bibitem{square-data}
R. H. Hardin, \emph{A068293}, The On-Line Encyclopedia of Integer Sequences,
February 24, 2002, NAME and first COMMENT;
identified with four-colour king boards in A212209, row 4.
\url{https://oeis.org/A068293}; \url{https://oeis.org/A212209}.
\end{thebibliography}
\noindent\small
\copyright\ 2026 Haobo Ma and Wenlin Zhang.
Licensed under Creative Commons Attribution 4.0 International
(CC BY 4.0), \url{https://creativecommons.org/licenses/by/4.0/}.
\end{document}
